\section{Where Truth Will Dwell}

\begin{quotex}
Possessing these new faculties, freed from excessive labour, the new man will have the power to organize life in the new society and, with the help of his clairvoyance, he will establish a new order where \emph{truth will dwell}.

\end{quotex}
In the brief monograph, \emph{The Problem of the New Man}, \textbf{Boris Mouravieff} describes the historical conditions leading up to the issues of today and the necessary solution that will usher in the next Golden Age, ending the current cycle. Mouravieff claims to be revealing the esoteric teachings of the Eastern Churches, which have not been made public until now.

\paragraph{Four moments of the Spirit}
He identifies four dominant trends, which we are calling the moments of the Spirit since they are the historical manifestation of a single idea. These are: Philosophy, Religion, Science, and Art.

\begin{description}
\item[Philosophy ] Philosophy is based on the concept which is both universal as idea and concrete as it is expressed as the judgment of particular states of affairs. This judgment is the unity of subject and object, or essence and existence. 

\item[Religion ] Religion denies that reality can be fully expressed conceptually. For the religious consciousness, the concept and the thing, essence and existence are no longer statically givens, as objects of contemplation. The ethical will connects the ideal to existence, which it actualizes. 

\item[Science ] Science, the child of nominalism, then reduces the concept to the practical. Bracketing out any reference to the transcendental, it is autonomous with respect to the ethical. The will is replaced with volition, or desire, whose aim is merely practical. Ideas, concepts, and theories are tools to manipulate the world for practical ends. 

\item[Art ] Art is knowledge of the individual, not the universal concept of philosophy. Its method is intuition, a direct, unmediated, unity. But intuition is real only to the extent to which it is expressed. For intuition, the ethical is the aesthetic. 

\end{description}
\paragraph{Historical Relationships}
Mouravieff does not begin his story with the unity of Spirit and virtues of the archaic Greek and Roman kingdoms, but with the philosophical era that came to dominate in the period of the Empire. Besides the four main streams of Platonism, Aristotelianism, Stoicism, and Epicureanism, there were many minor movements that sprang up such as Cynicism, Skepticism and so on. Plotinus marks the end of philosophical moment and his students Boethius and Augustine mark the transition to the religious moment.

After the ideal of the contemplative thinker, the knight became the new elite. The knight was distinguished by his physical strength and martial prowess\footnote{\url{https://gornahoor.net/?p=277}}. The knight acted according to the ethical will of Christendom.

Military technology and the rise of nominalism replaced the power of the knight and the authority of the Church. The ``intellectual" class became the new elite. Exploration, discovery, and economic activity began to dominate. No longer bound by the ethical precepts of the Medieval period, money become the new power. Thus an economic and ideological superstructure was imposed over brute reality, becoming in its place a sort of ``second reality". The manipulation of this second reality is now the dominant interest of the elite, despite the many irruptions of first reality that upset its most ambitious schemes.

In our time, we see that money has changed from the medium of exchange to a commodity itself. It is no longer the industrialist who generates wealth, but rather the financier who knows how to manipulate money instruments. This is almost a third reality, hence it cannot endure in its present from. The question is whether this collapse will be gentle or chaotic, and whether something higher can replace it.

\paragraph{The Age of the Spirit}
Art includes the three preceding moments without deforming them. Just as the gadgets of the technocrat of today would appear as the product of supermen to a Knight, so the coming age of the Spirit will require a new elite of supermen. The principle faculty of the elite man to come ``will be the capacity to distinguish the real from the false, the truth from lies, spontaneously, without evidence or other support," that is, a direct intuitive understanding. The problem today is that man is unable to distinguish truth from lies; all his leaders lie to him, whether deliberately or \emph{de facto} since they themselves don't even know what is true.

\paragraph{A New Beginning}
\begin{quotex}
Looking for and hasting unto the coming of the day of the Lord, by which the heavens being on fire shall be dissolved, and the elements shall melt with the burning heat? But we look for new heavens and a new earth according to his promises, in which justice dwelleth. \flright{\textsc{2 Peter 3:12-13}}

\end{quotex}
Just as the Knight dominated by muscular power and the technocrat by intellectual power, the new man will lead through the power of Spirit, which will give him new abilities. First of all, Mouravieff makes clear that Jesus made the distinction between his disciples and the ``world", those on the ``outside". He prayed to God not for all men, but for those ``whom thou hast given me." The incarnation of the Logos is a process that unfolds over time. It is the task of the disciples to facilitate that unfolding through their own efforts. Mouravieff writes:

\begin{quotex}
Between hope and realization there is still a path to travel. This distance must be covered by ourselves, by our own efforts. Nothing is given free. Everything must be paid for.

\end{quotex}
Interestingly, Mouravieff claims that the natural man will continue to lie as long as his thoughts are private. Hence, the ability to read minds will be a faculty of the superman. To gain these faculties, the new elite must be ``twice born". While the current age seeks the transformation of nature through technology, the new age will transform man through a Positivism of the Spirit\footnote{\url{https://gornahoor.net/?p=1189}}.

\paragraph{The Projection of the I}
Mouravieff writes that the Eastern Orthodox esoteric tradition 

\begin{quotex}
teaches that all civilization is only the projection on the outside world of the consciousness of the I of the elite man.

\end{quotex}
Thus, the elite of any era are able to impose their world view on the masses. In the past, this has not been the real I, but ``unstable, composite and multiple" counterfeits, in conflict with each other. The new elite man will have consciousness of the ``real I". What had been privation previously will now come under the dominion of this elite.

\begin{quotex}
Thus the edifice of the future civilization will no longer be built by him on ``sand", but on the ``rock" of his self-consciousness, on the divine spark. 

\end{quotex}
References:

\emph{The Problem of the New Man}, Boris Mouravieff

\emph{Western Philosophy} Vol 5, A. Robert Caponigri

Originally appeared in the Meditations on the Tarot blog on 9 April 2011



\flrightit{Posted on 2023-02-07 by Cologero }

\begin{center}* * *\end{center}

\begin{footnotesize}\begin{sffamily}



\texttt{Ismo Meinander on 2011-04-09 at 05:58 said: }

This sounds like `a true new age teaching'.

If I remember correctly, Evola spoke about this coming age and the real superman – or the New Man – in `Ride the Tiger' and in reference to the `third age of the Spirit' which is a conception made by Joachim de Flore. I think that in these weird transitional times there are many signs of this happening and the `new age fallacies' are a somewhat naive and many times dangerous symptoms of this: new faculties are developing and they are in many cases spontaneous, and in the general climate of our times they lead especially young people without effective means of understanding and protection of tradition to madness; since the new faculties are not given proper spiritual orientation they lead up easily to the most distressing developments of psychism.


\hfill


\end{sffamily}\end{footnotesize}
